\documentclass{beamer}
\usepackage{ctex, hyperref}
\usepackage[T1]{fontenc}%使用8位编码的T1字体编码加载字体编码

% other packages
\usepackage{latexsym,amsmath,xcolor,multicol,booktabs,calligra}
\usepackage{graphicx,pstricks,listings,stackengine}
% 列表缩进与间距：不要用 enumitem，直接在 beamer 中手动设定
\setlength{\leftmargini}{1.4em}
\setlength{\leftmarginii}{1.2em}

\setbeamerfont{block title}{size=\normalsize}
\setbeamerfont{block body}{size=\small}
\setbeamerfont{alertblock title}{size=\normalsize}
\setbeamerfont{alertblock body}{size=\small}
\setlength{\leftmargini}{1.4em}
\setlength{\leftmarginii}{1.2em}


\author[Zhennan Zhao]{
{\zihao{-4}\hspace{0.1em}赵\hspace{0.32em}圳\hspace{0.32em}楠\hspace{0.1em}}\\[2pt]
{\scalebox{0.58}[0.58]{Zhennan Zhao}}
}
\title{2024级直博生综合考试}
\subtitle{个人情况、科研准备与研究设想}
\institute{上海大学理学院数学系}
\date{2026年3月12日}
\usepackage{SHU}

% defs
\def\cmd#1{\texttt{\color{red}\footnotesize $\backslash$#1}}
\def\env#1{\texttt{\color{blue}\footnotesize #1}}
\definecolor{deepblue}{rgb}{0,0,0.5}
\definecolor{deepred}{rgb}{0.5,0,0}
\definecolor{deepgreen}{rgb}{0,0.5,0}
\definecolor{halfgray}{gray}{0.55}

\lstset{
    basicstyle=\ttfamily\small,
    keywordstyle=\bfseries\color{deepblue},
    emphstyle=\ttfamily\color{deepred},    % Custom highlighting style
    stringstyle=\color{deepgreen},
    numbers=left,
    numberstyle=\small\color{halfgray},
    rulesepcolor=\color{red!20!green!20!blue!20},
    frame=shadowbox,
}



















% ========================================文件开始========================================
\begin{document}
\kaishu

% ========================================
% 封面
% ========================================
\begin{frame}
    \titlepage
\end{frame}

% ========================================
% 目录
% ========================================
\begin{frame}{汇报提纲}
    \tableofcontents[sectionstyle=show,subsectionstyle=hide]
\end{frame}
 
% ========================================
\section{个人简介与研究兴趣}
% ========================================
\begin{frame}{个人简介}
\small

\begin{columns}[T,totalwidth=\textwidth]
    \begin{column}{0.57\textwidth}
        {\setlength{\itemsep}{0.45em}
        \begin{itemize}
            \item \textbf{姓名：}赵圳楠
            \item \textbf{培养层次：}2024级数学学科直博生
            \item \textbf{所属单位：}上海大学理学院数学系
            \item \textbf{导师：}赖耕教授
            \item \textbf{拟研究方向：}偏微分方程（PDE）
            \item \textbf{重点关注：}非线性双曲守恒律与黎曼问题
        \end{itemize}
        }
    \end{column}


    \begin{column}{0.39\textwidth}
        \begin{alertblock}{科研兴趣关键词}
            {\setlength{\itemsep}{0.40em}
            \begin{itemize}
                \item 双曲守恒律
                \item 可压 Euler 方程组
                \item 一维/二维黎曼问题
                \item 激波、稀疏波、接触间断
                \item 波的相互作用与稳定性
            \end{itemize}
            }
        \end{alertblock}
    \end{column}
\end{columns}

\vspace{0.35em}

\begin{block}{个人定位}
现阶段定位为：以\textbf{双曲型偏微分方程}为核心方向，
以\textbf{黎曼问题}为主要切入点，
在导师指导下逐步学习守恒律理论、Euler 方程组及典型波现象的分析方法，
为后续开展实质性研究做好准备。
\end{block}

\end{frame}

\begin{frame}{汇报思路：以科研准备为主线}
\begin{block}{本次汇报重点}
相较于一般性的个人介绍，本次汇报更侧重展示本人在\textbf{科研方向选择、知识储备、问题理解和后续研究计划}方面的准备情况。
\end{block}

\begin{itemize}
    \item 已完成课程学习，具备继续开展博士研究的基础条件；
    \item 研究兴趣与导师方向高度契合，聚焦\textbf{非线性双曲守恒律方程组及其黎曼问题}；
    \item 当前重点是夯实理论基础、系统阅读文献、凝练可操作的研究问题；
    \item 后续希望在双曲型PDE若干典型模型上形成稳定的研究积累。
\end{itemize}
\end{frame}




% ========================================
\section{课程基础与学术准备}
% ========================================



% ========================================
% 课程基础与学术准备：微调美化版
% ========================================

% ----------------------------------------
% 第1页：课程修读概览
% 调整点：
% 1. 指标字体缩小
% 2. 下半页补充“课程完成情况解读”
% 3. 页面上下信息更均衡
% ----------------------------------------
\begin{frame}{课程修读概览}
\vspace{-1.2em} % 【修正1】大幅上移整体内容，充分利用标题下方空闲的白色区域

% ================= 上半部分：四个数据框 =================
\begin{columns}[T, onlytextwidth]
    \begin{column}{0.48\textwidth}
        \begin{block}{已修总学分}
            % 【修正2】将上半部分方框高度压缩至 0.5cm
            \begin{minipage}[c][0.5cm][c]{\textwidth}
                \centering
                {\small 57.0}
            \end{minipage}
        \end{block}
        
        \vspace{0.05em} % 极限压缩上下框间距
        
        \begin{block}{加权平均分}
            \begin{minipage}[c][0.5cm][c]{\textwidth}
                \centering
                {\small 90.63}
            \end{minipage}
        \end{block}
    \end{column}
    
    \begin{column}{0.48\textwidth}
        \begin{block}{课程通过情况}
            \begin{minipage}[c][0.5cm][c]{\textwidth}
                \centering
                {\small 全部通过}
            \end{minipage}
        \end{block}
        
        \vspace{0.05em}
        
        \begin{block}{加权绩点}
            \begin{minipage}[c][0.5cm][c]{\textwidth}
                \centering
                {\small 3.77 / 4.0}
            \end{minipage}
        \end{block}
    \end{column}
\end{columns}

\vspace{0.1em} % 上下两部分的过渡间距收紧

% ================= 下半部分：两个长文本框 =================
\begin{columns}[T, onlytextwidth]
    \begin{column}{0.48\textwidth}
        \begin{block}{课程完成情况解读}
            % 【修正3】高度从 3.4cm 大幅缩减至 2.8cm，确保方框底边露出水面！
            \begin{minipage}[t][2.8cm][t]{\textwidth}
                \footnotesize
                \vspace{0.1em}
                \begin{itemize}
                    \setlength{\itemsep}{0.1em} % 同步压缩列表间距，防止文字溢出方框
                    \item 已较系统完成博士阶段前期课程修读任务；
                    \item 成绩整体稳定，说明课程学习状态较为扎实；
                    \item 当前已具备从课程训练逐步过渡到研究训练的基础条件。
                \end{itemize}
            \end{minipage}
        \end{block}
    \end{column}
    
    \begin{column}{0.48\textwidth}
        \begin{alertblock}{阶段性判断}
            % 与左侧保持绝对一致的 2.8cm 高度
            \begin{minipage}[t][2.8cm][t]{\textwidth}
                \footnotesize
                \vspace{0.1em}
                目前已完成现阶段培养方案要求的主要课程训练。
                
                \vspace{0.5em}
                下一阶段的重点将由“课程学习”逐步转向“文献阅读、问题理解与研究准备”。
            \end{minipage}
        \end{alertblock}
    \end{column}
\end{columns}

\end{frame}

% ----------------------------------------
% 第2页：课程结构与方向关联
% 调整点：
% 1. 三列模块更紧凑
% 2. 减少零散断行
% 3. 底部总结更自然
% ----------------------------------------
\begin{frame}{课程结构与方向关联}
\small

\begin{columns}[T,totalwidth=\textwidth]
    \begin{column}{0.32\textwidth}
        \begin{block}{分析基础}
        \footnotesize
        \textbf{代表课程：}\\
        泛函分析（98）、实分析（84）、拓扑学（96）、几何分析（92）

        \vspace{0.45em}
        \textbf{主要作用：}\\
        为函数空间、收敛与估计、抽象结构理解提供基础。
        \end{block}
    \end{column}

    \begin{column}{0.32\textwidth}
        \begin{block}{PDE 相关}
        \footnotesize
        \textbf{代表课程：}\\
        微分方程（86）、可积偏微分方程（97）、物理学与偏微分方程（78）

        \vspace{0.45em}
        \textbf{主要作用：}\\
        帮助理解 PDE 模型、解的结构及波现象的初步认识。
        \end{block}
    \end{column}

    \begin{column}{0.32\textwidth}
        \begin{block}{方法补充}
        \footnotesize
        \textbf{代表课程：}\\
        有限元方法（96）、统计学习与人工智能基础（89）、学术英语写作与交流（94）

        \vspace{0.45em}
        \textbf{主要作用：}\\
        补充数值视角、交叉方法与学术表达能力。
        \end{block}
    \end{column}
\end{columns}

\vspace{0.3em}

\begin{alertblock}{与当前研究方向直接相关的课程}
\footnotesize
泛函分析、微分方程、可积偏微分方程、有限元方法，
构成了我进一步进入\textbf{双曲型偏微分方程、守恒律与黎曼问题}研究的主要课程基础。
\end{alertblock}

\end{frame}

% ----------------------------------------
% 第3页：课程训练对后续研究的支撑
% 调整点：
% 1. 压缩 2×2 block 的文字密度
% 2. 结论整体上移
% 3. 用 alertblock 承接结论，不再让底部显得松散
% ----------------------------------------
\begin{frame}{课程训练对后续研究的支撑}
\small

\setlength{\columnsep}{0.6em}

\begin{columns}[T,totalwidth=\textwidth]
    \begin{column}{0.49\textwidth}
        \begin{block}{理论准备}
        \begin{minipage}[t][1.55cm][t]{\linewidth}
        \footnotesize
        \raggedright
        通过泛函分析、实分析、拓扑学等课程训练，
        我逐步建立了对函数空间、收敛性、先验估计与严格证明方式的理解。
        \end{minipage}
        \end{block}

        \vspace{0.02em}

        \begin{block}{方法补充}
        \begin{minipage}[t][1.35cm][t]{\linewidth}
        \footnotesize
        \raggedright
        有限元方法等课程补充了数值分析视角，
        有助于从计算与现象层面理解 PDE 模型及其解的行为。
        \end{minipage}
        \end{block}
    \end{column}

    \begin{column}{0.49\textwidth}
        \begin{block}{方向进入}
        \begin{minipage}[t][1.55cm][t]{\linewidth}
        \footnotesize
        \raggedright
        微分方程、可积偏微分方程等课程，
        帮助我逐步建立对 PDE 基本结构、典型波形和研究问题的认识。
        \end{minipage}
        \end{block}

        \vspace{0.02em}

        \begin{block}{学术训练}
        \begin{minipage}[t][1.35cm][t]{\linewidth}
        \footnotesize
        \raggedright
        学术英语写作与交流、学术规范与写作等课程，
        为后续文献阅读、学术表达与研究写作提供了基础准备。
        \end{minipage}
        \end{block}
    \end{column}
\end{columns}

\vspace{0.06em}

\begin{alertblock}{本部分结论}
\footnotesize
课程学习的价值不只是“完成成绩要求”，更重要的是帮助我完成了从
“系统学习知识”到“逐步进入研究问题”的过渡。
\end{alertblock}

\end{frame}















% ========================================
\section{科研经历、学习方向与内容准备}
% ========================================
\begin{frame}{科研经历与现阶段学术积累}
\begin{itemize}
    \item 当前处于博士培养早期阶段，暂未形成正式论文成果；
    \item 现阶段工作的重点是：\textbf{课程学习、方向聚焦、文献阅读、问题准备}；
    \item 已将兴趣集中于\textbf{非线性双曲守恒律方程组及其黎曼问题}，并以导师研究方向为参照进行系统性学习规划；
    \item 在竞赛与学生工作中积累了建模、表达和组织协调能力：
    \begin{itemize}
        \item 华为杯研究生数学建模竞赛三等奖；
        \item 担任团支书，负责组织相关活动。
    \end{itemize}
\end{itemize}

\begin{alertblock}{本页表述策略}
本次PPT更突出“\textbf{科研准备和研究潜力}”，即通过系统学习与问题意识展示进入博士研究阶段的可持续性。
\end{alertblock}
\end{frame}

\begin{frame}{围绕导师方向拟系统学习的三个模块}
\begin{block}{模块一：双曲守恒律与一维黎曼问题}
重点学习守恒律基本理论、弱解与熵条件、激波与稀疏波结构、波前跟踪思想，以及典型一维黎曼问题的求解框架。
\end{block}

\begin{block}{模块二：可压Euler方程组与二维黎曼问题}
重点学习可压流体模型、二维自相似解结构、四象限型初值问题、接触间断与激波相互作用，以及二维几何结构带来的新困难。
\end{block}

\begin{block}{模块三：复杂流动中的自由边界与反应流问题}
结合导师方向，重点关注激波反射、超声速射流、爆燃/爆轰波、气体向真空扩散等问题中自由边界、波系耦合与稳定性分析。
\end{block}
\end{frame}

\begin{frame}{已纳入PPT展示的“科研学习内容”}
\small
\begin{enumerate}
    \item \textbf{基础理论层面：}
    已系统梳理双曲型方程、守恒律、弱解、熵解、特征线方法等核心概念；
    \item \textbf{模型理解层面：}
    重点关注可压Euler方程组及其在激波、稀疏波、接触间断中的基本波结构；
    \item \textbf{问题意识层面：}
    关注一维到二维推广中出现的结构复杂性，以及波相互作用、自由边界、真空状态等难点；
    \item \textbf{方法层面：}
    初步了解特征分解、先验估计、自相似变换、构造解与稳定性分析等研究思路；
    \item \textbf{文献阅读层面：}
    围绕二维黎曼问题、激波反射、超声速射流等主题开展阶段性阅读和笔记整理。
\end{enumerate}

\begin{alertblock}{说明}
本页内容适合作为“\textbf{当前已开展的方向性学习}”进行展示，后续可根据你的实际阅读进度继续细化或替换。
\end{alertblock}
\end{frame}

% ========================================
\section{研究方向理解与问题认识}
% ========================================
\begin{frame}{对导师研究方向的理解}
\begin{itemize}
    \item 赖耕老师的研究主要集中在\textbf{非线性双曲守恒律方程组的理论及应用}；
    \item 具体包括\textbf{可压Euler方程组的二维黎曼问题}，以及气体向真空扩散、激波反射、管道流、超声速射流、爆燃波和爆轰波等数学问题；
    \item 这一方向兼具\textbf{理论深度}与\textbf{应用背景}：一方面关注解的存在性、结构与稳定性，另一方面与流体、燃烧和高速气体动力学问题密切相关；
    \item 对我而言，该方向的吸引力在于：既有明确的数学问题导向，也有丰富的物理背景支撑，适合作为博士阶段长期深入的研究主线。
\end{itemize}
\end{frame}

\begin{frame}{我对该方向核心问题的理解}
\begin{block}{为什么选择黎曼问题作为切入点？}
黎曼问题是双曲守恒律研究中的基本模型，它能够集中体现激波、稀疏波、接触间断等典型波结构的形成机制，是理解更复杂初边值问题的重要起点。
\end{block}

\begin{itemize}
    \item \textbf{从一维到二维：}二维问题中的波相互作用更复杂，几何结构和自由边界现象更加突出；
    \item \textbf{从理想到实际：}真空、反应项、边界效应、喷流和反射等因素会显著提高问题难度；
    \item \textbf{从定性到定量：}不仅需要知道“会出现什么波”，还需要分析其相互作用机制、存在区间和稳定性特征。
\end{itemize}

\begin{alertblock}{个人认识}
我希望以黎曼问题为基本切入点，在扎实掌握经典理论的基础上，逐步过渡到二维、自由边界和反应流等更复杂的问题。
\end{alertblock}
\end{frame}

% ========================================
\section{博士阶段研究设想}
% ========================================
\begin{frame}{博士阶段拟开展研究工作的设想}
\begin{enumerate}
    \item \textbf{第一阶段：夯实理论基础}
    \begin{itemize}
        \item 深入学习双曲守恒律、Euler方程组、弱解与熵解理论；
        \item 系统阅读导师团队相关工作与经典文献，形成自己的知识框架。
    \end{itemize}

    \item \textbf{第二阶段：凝练研究问题}
    \begin{itemize}
        \item 在导师指导下聚焦一个具体模型或具体场景；
        \item 从二维黎曼问题、激波反射、真空边界或超声速射流等方向中寻找可切入问题。
    \end{itemize}

    \item \textbf{第三阶段：开展实质研究}
    \begin{itemize}
        \item 尝试从存在性、唯一性、结构分析、稳定性等角度开展工作；
        \item 在必要时结合数值实验辅助理解理论现象，争取形成阶段性成果。
    \end{itemize}
\end{enumerate}
\end{frame}

\begin{frame}{当前更具体的学习/研究切入设想}
\small

\setlength{\columnsep}{0.6em}

\begin{columns}[T,totalwidth=\textwidth]
    \begin{column}{0.32\textwidth}
        \begin{block}{方向A：黎曼问题}
        \begin{minipage}[t][3.15cm][t]{\linewidth}
        \footnotesize
        \raggedright
        {\setlength{\itemsep}{0.38em}
        \begin{itemize}
            \item 自相似框架下的波结构分类
            \item 接触间断与激波的相互作用
            \item 特定初值构型下解的局部与整体行为
        \end{itemize}}
        \end{minipage}
        \end{block}
    \end{column}

    \begin{column}{0.32\textwidth}
        \begin{block}{方向B：自由边界}
        \begin{minipage}[t][3.15cm][t]{\linewidth}
        \footnotesize
        \raggedright
        {\setlength{\itemsep}{0.38em}
        \begin{itemize}
            \item 真空边界附近的解结构
            \item 超声速喷流中的拦截激波
            \item 几何边界对流动形态的影响
        \end{itemize}}
        \end{minipage}
        \end{block}
    \end{column}

    \begin{column}{0.32\textwidth}
        \begin{block}{方向C：反应流}
        \begin{minipage}[t][3.15cm][t]{\linewidth}
        \footnotesize
        \raggedright
        {\setlength{\itemsep}{0.38em}
        \begin{itemize}
            \item 反应项引入后的波系耦合
            \item 爆燃/爆轰波中的传播机制
            \item 数学模型与物理背景的对应关系
        \end{itemize}}
        \end{minipage}
        \end{block}
    \end{column}
\end{columns}

\vspace{0.1em}

\begin{alertblock}{总体原则}
\footnotesize
以\textbf{黎曼问题}为主要切入，优先从\textbf{可控、经典、可积累}的问题入手，
再逐步过渡到\textbf{自由边界、喷流与反应流}等更复杂模型；
在推进过程中同步关注\textbf{数学结构、物理背景与分析方法}之间的对应关系。
\end{alertblock}

\end{frame}

% ========================================
\section{总结}
% ========================================
\begin{frame}{总结与展望}
\begin{itemize}
    \item 已完成现阶段课程学习，具备较扎实的分析基础与PDE基础；
    \item 研究兴趣与导师方向高度契合，已将重点聚焦于\textbf{非线性双曲守恒律与黎曼问题}；
    \item 现阶段更重视\textbf{科研学习、问题意识与方向准备}，而不仅是简单的成果罗列；
    \item 后续希望在导师指导下，逐步形成稳定的研究主题、研究方法与阶段性成果。
\end{itemize}

\vspace{0.8em}
\begin{center}
{\Large 恳请各位老师批评指正！}
\end{center}
\end{frame}

\begin{frame}
    \begin{center}
        {\Huge Thanks!}
    \end{center}
\end{frame}

\end{document}
